移动端 GUI 智能体有时必须输入它在许多屏之前看到过的文字。我们把这类步骤称为\textbf{延迟信息需求}：它在 AITZ 上一例也没有（18,643 步中 0 例），在 GUI-Odyssey 的抽样轨迹中则有 31\% 至少含一例，源屏与需求步之间的距离中位数为 5 屏。要挖掘这类需求必须依赖一个验证器：所输入的文字出现在某个更早屏的 OCR（光学字符识别）文本中，且不出现在当前屏上。我们指出，同一个验证器顺带免费给出了训练标签——它把每个更早屏标记为"含有该文字"或"不含"，而正是这种免标注监督让回看真正奏效。用这些标签训练后，全历史读取器在 188 个已验证样例上的源屏定位率从 prompt 下的 54.8\% 升到 65.4\%，逐屏打分器从 50.0\% 升到 72.3\%。两个训练变体之间定位相差 6.9 个点（\pval{0.086}）、端到端相差 2.1 个点（\pval{0.45}）。打分器的优势集中在 8--12 屏的距离区间（78.2\% 对 56.4\%，\pval{0.003}），那里一次调用需要同时权衡许多候选。在一批与所有设计决策完全隔离的第四轮抽样轨迹上，打分器定位 79.6\%、训练后的全历史读取器 76.3\%（\pval{0.63}）、prompt 全历史 59.1\%（\pval{0.003}）；带历史读取把端到端准确率从 prompt 全历史的 23.7\% 提到 49.5\% 与 51.6\%（两者均 $p<0.001$）。我们同时报告两个训练变体共同失效的区间，以及任何研究该需求都必须施加的两个评测过滤条件。
